\chapter{结论}\label{chap:conclusion}

\section{与现有工作的比较}\label{sec:comparison_related}

本文的工作与现有PINN期权定价研究的主要区别体现在以下几个维度。

在多模型统一性方面，现有工作（Dhiman等\cite{dhiman2023pinn}、Bae等\cite{bae2024localvol}、Hainaut等\cite{hainaut2024heston}）均针对单一模型，本文是首个将BSM、CEV、Heston三大模型统一在单一网络中的工作。多模型联合训练带来了互相增益：BSM/CEV的简单结构为Heston提供了良好的初始化，有助于克服Heston二维PDE的收敛困难。与参数化PINN相比，统一PINN在BSM和CEV上精度更高，在Heston上参数化PINN略优，但统一PINN的网络参数量仅为参数化PINN的约四分之一，具有更高的参数效率，说明软掩码和加法参数化实现了更高效的归纳偏置。

在交互方式方面，现有PINN工具均为纯数值接口，用户需手动指定模型类型和参数。本文引入LLM路由层，支持自然语言输入，是PINN与LLM结合用于期权定价的首次尝试。与PDE-AGENT\cite{liu2025pdeagent}相比，本文的LLM仅作为路由层，不参与求解过程，因此延迟更低、可靠性更高。

在约束处理方面，本文采用加法参数化输出，以Black-Scholes解析解为基线，这与Zamxaka\cite{zamxaka2023cev}的参数化PINN思路相近，但本文将其推广至多模型统一求解，并结合软掩码实现了模型间的连续插值。

\section{工作总结}\label{sec:summary}

本文研究了基于PINN的多模型期权定价方法，通过软掩码机制将BSM、CEV、Heston三大模型的PDE算子统一为单一参数化形式，以加法参数化输出解决深度虚值区域的数值退化问题，以LLM路由层支持自然语言输入。实验结果表明，三种模型的ATM相对误差均低于1\%，推断速度约2ms，LLM路由层准确率100\%，消融实验验证了各核心组件的必要性。

\section{局限性}\label{sec:limitations}

本文的工作存在以下局限性，需要在未来工作中进一步改进。

首先，PINN-Heston在真实市场数据上的泛化能力不足：v16的Heston训练参数范围（$\kappa\in[1.0,3.0]$，$\xi\in[0.2,0.4]$，$\rho\in[-0.7,-0.5]$）无法覆盖SPY期权校准得到的极端参数（$\kappa$可达20，$\xi$可高达2.1，$\rho$可达$-0.67$），且统一PINN用同一个6层128宽的网络同时拟合三种模型，Heston的5维参数空间对网络容量要求远高于BSM/CEV，构成根本性的容量瓶颈。在标准参数集下，统一PINN的Heston精度良好，但扩展至真实市场的宽参数范围仍是重要的未来工作。此外，当前方法仅针对欧式看涨期权，不支持美式期权（需要处理自由边界问题）或奇异期权（障碍期权、亚式期权等）。

另一方面，LLM路由层依赖DeepSeek API，需要网络连接，且存在约500ms至2000ms的API调用延迟，未来可以考虑使用本地部署的小型LLM（如Qwen-7B）替代。当前框架仅输出点估计，不提供预测不确定性，而在金融应用中，不确定性量化（如置信区间）对风险管理至关重要，未来可以引入贝叶斯PINN或集成方法提供不确定性估计。

\section{未来工作}\label{sec:future}

基于上述局限性，本文提出以下未来工作方向。

A股市场适配方面，当前框架的训练参数范围针对美股市场校准，A股市场（上证50ETF期权、沪深300ETF期权）的隐含波动率特征与美股有显著差异：A股波动率通常更高（年化波动率常超过30\%），且受政策因素影响，均值回归速度$\kappa$和长期均值$\theta$的分布与美股不同。未来可针对A股市场重新校准参数范围，并适配AkShare等国内数据接口，同时处理涨跌停限制、T+1交割等市场微结构特殊规则对定价的影响。

美式期权扩展方面，将本文方法扩展至美式期权，处理自由边界（最优执行边界）问题。美式期权的定价等价于最优停止问题，对应的PDE为变分不等式，需要引入惩罚函数法或互补条件。可以借鉴现有的PINN美式期权工作\cite{louskos2021american}，将自由边界条件纳入统一损失函数。奇异期权方面，障碍期权（Barrier Options）的定价PDE与欧式期权相同，但边界条件不同（在障碍处$V=0$或$V=\text{回扣}$），理论上可纳入本文框架，需额外设计障碍边界的网络处理方式；亚式期权（Asian Options）的收益依赖路径平均价格，引入路径依赖变量后维度更高，需要更大规模的网络和采样策略。

Greeks计算方面，利用PINN的自动微分能力，直接计算期权的Greeks（Delta、Gamma、Vega等），无需有限差分近似，这是PINN相对于传统数值方法的天然优势，在风险管理中具有重要价值。本地LLM部署方面，将LLM路由层替换为本地部署的小型语言模型，降低API依赖和调用延迟，同时保护用户的期权参数隐私。

\section{本章小结}\label{sec:conclusion_summary}

本文研究了基于PINN的多模型期权定价方法，通过软掩码机制、加法参数化输出和LLM路由层，实现了BSM、CEV、Heston三大模型的统一求解和自然语言交互。实验结果表明，本文方法在精度和效率上均达到实用水平，为PINN在金融工程中的应用提供了新的思路，也为LLM与科学计算工具的结合提供了一个具体的实践案例。
